\section*{Capítulo \ref{ch:b3:bioinorganic} --- Química bioinorgânica}

\begin{solution}{exo:b3:bioinorganic:1}
$\log(\theta/(1 - \theta))$ vai de $\log 0.25 = -0.602$ a $\log 4 = +0.602$ enquanto $\log p$ aumenta
$\log(5.9/2.1) = 0.449$: $n = 1.204/0.449 = 2.7$.
\end{solution}

\begin{solution}{exo:b3:bioinorganic:2}
$\ce{Mn^2+} < \ce{Fe^2+} < \ce{Co^2+} < \ce{Ni^2+} < \ce{Cu^2+} > \ce{Zn^2+}$. O zinco ($d^{10}$) não tem estabilização do campo ligante nem
distorção de Jahn--Teller: ele volta a ficar abaixo do cobre.
\end{solution}

\begin{solution}{exo:b3:bioinorganic:3}
$\ce{Ca^2+}$ e $\ce{Fe^3+}$, duros: carboxilatos (e fenolatos para o ferro). $\ce{Cu+}$ e $\ce{Hg^2+}$, moles:
tiolatos de cisteína. $\ce{Zn^2+}$, intermediário: histidina e cisteína.
\end{solution}

\begin{solution}{exo:b3:bioinorganic:4}
$\ce{2H2O -> O2 + 4H+ + 4e-}$: quatro elétrons, quatro fótons por $\ce{O2}$; quatro mols de fótons por
mol de $\ce{O2}$.
\end{solution}

\begin{solution}{exo:b3:bioinorganic:5}
Mioglobina: $\theta = 13/13.37 = 0.972$ e $5/5.37 = 0.931$: ela cede $(0.972 - 0.931)/0.972 = 4$\,\% de seu
oxigênio. Hemoglobina: $\theta = 0.972$ e $0.724$: ela cede 26\,\%.
\end{solution}

\begin{solution}{exo:b3:bioinorganic:6}
$[\mathrm{HbCO}]/[\mathrm{HbO_2}] = 230 \times \num{50e-6}/0.2095 = 0.055$; fração $0.055/1.055 = 5.2$\,\%.
\end{solution}

\begin{solution}{exo:b3:bioinorganic:7}
Sulfetos $-2$ cada, cisteinatos $-1$ cada. $\ce{[Fe2S2(SR)4]^2-}$: soma dos ferros $-2 + 4 + 4 = +6$: dois
Fe(III). $\ce{[Fe2S2(SR)4]^3-}$: $+5$, um Fe(III) e um Fe(II) (valência mista). $\ce{[Fe4S4(SR)4]^2-}$:
$-2 + 8 + 4 = +10$, dois Fe(III) e dois Fe(II), cada ferro $+2.5$ em média.
\end{solution}

\begin{solution}{exo:b3:bioinorganic:8}
$\ce{[PtCl2(NH3)2] + H2O -> [PtCl(H2O)(NH3)2]+ + Cl-}$. A alta concentração de cloreto do sangue
desloca o equilíbrio de volta; dentro das células, o cloreto é muito menor, e o
aquacomplexo se forma. A platina se liga ao N7 da guanina. No isômero \textit{trans}, os dois
sítios \omterm{def:b3:complex-mechanisms:labile}{lábeis} são opostos e não conseguem alcançar duas bases adjacentes de uma mesma
fita.
\end{solution}

\begin{solution}{exo:b3:bioinorganic:9}
$1/T_1 = 1/\qty{1.2}{s} + \qty{4.0}{L.mmol^{-1}.s^{-1}} \times \qty{0.10}{mmol/L} = 0.833 + 0.400 = \qty{1.23}{s^{-1}}$: $T_1 = \qty{0.81}{s}$.
\end{solution}

\begin{solution}{exo:b3:bioinorganic:10}
$K = K_{\mathrm{PbL}}/K_{\mathrm{CaL}} = 10^{18.0 - 10.7} = 10^{7.3} = \num{2e7}$: o chumbo desloca completamente o cálcio. O ligante
livre ligaria também o cálcio do sangue e abaixaria perigosamente sua concentração;
o complexo de cálcio troca com o chumbo e deixa inalterado o nível
de cálcio.
\end{solution}

\begin{solution}{exo:b3:bioinorganic:11}
Muito abaixo de $K_M$, a reação é de primeira ordem: $k' = (k_{\mathrm{cat}}/K_M)[\mathrm E]$, $t_{1/2} = \ln 2/k'$. Anidrase
carbônica: $k' = \qty{100}{s^{-1}}$, $t_{1/2} = \qty{6.9}{ms}$. \omterm{def:b3:complex-kinetics:enzyme}{Enzima} típica: $k' = \qty{0.10}{s^{-1}}$, $t_{1/2} = \qty{6.9}{s}$, mil
vezes mais, lento demais para o segundo que uma hemácia passa em um capilar pulmonar.
\end{solution}

\begin{solution}{exo:b3:bioinorganic:12}
Metade dos sítios com CO significa $[\mathrm{HbCO}] = [\mathrm{HbO_2}]$: $p(\ce{CO}) = p(\ce{O2})/M = 0.2095/230 = \qty{9.1e-4}{bar}$, cerca de
\qty{910}{ppm}.
\end{solution}

\begin{solution}{pb:b3:bioinorganic:1}
\textbf{1.} $\theta = 13^{2.7}/(3.5^{2.7} + 13^{2.7}) = 0.972$.

\textbf{2.} $\theta(\qty{5.0}{kPa}) = 0.724$.

\textbf{3.} $5.0/(0.37 + 5.0) = 0.931$.

\textbf{4.} Seu $p_{50}$ é dez vezes menor: às pressões de um músculo em atividade, ela está muito mais
saturada que a hemoglobina à mesma pressão, de modo que o oxigênio passa da
hemoglobina para a mioglobina.

\textbf{5.} $n = 1$.

\textbf{6.} A ligação a um \omterm{def:b3:bioinorganic:haem}{heme} leva seu ferro para o plano e puxa a histidina e
sua hélice; as subunidades passam juntas para uma forma de maior afinidade nos
outros sítios.

\textbf{7.} $\qty{150}{g/L}/\qty{64500}{g/mol} = \qty{2.33e-3}{mol/L}$.

\textbf{8.} $4 \times \num{2.33e-3} = \qty{9.30e-3}{mol/L}$.

\textbf{9.} $\qty{9.30}{mmol/L} \times 0.972 = \qty{9.04}{mmol}$.

\textbf{10.} $\qty{9.30}{mmol/L} \times (0.972 - 0.724) = \qty{2.31}{mmol}$.

\textbf{11.} $\qty{2.31e-3}{mol} \times \qty{22.41}{L/mol} = \qty{52}{mL}$.

\textbf{12.} $0.248/0.972 = 26$\,\%.

\textbf{13.} $5.0^{2.7}/(4.0^{2.7} + 5.0^{2.7}) = 0.646$.

\textbf{14.} $\qty{9.30}{mmol/L} \times (0.972 - 0.646) = \qty{3.03}{mmol}$.

\textbf{15.} $3.03/2.31 = 1.31$: trinta por cento a mais.

\textbf{16.} O dióxido de carbono produzido pela respiração, hidratado a ácido carbônico pela
anidrase carbônica, e o ácido lático no esforço intenso.

\textbf{17.} $\qty{3.03}{mmol/L} \times \qty{5.0}{L/min} = \qty{15.1}{mmol/min}$.

\textbf{18.} $\qty{15.1e-3}{mol/min} \times \qty{22.41}{L/mol} = \qty{0.34}{L/min}$.

\textbf{19.} $230 \times \num{100e-6}/0.2095 = 0.110$.

\textbf{20.} $0.110/1.110 = 0.099$: cerca de 10\,\% dos sítios.

\textbf{21.} Os sítios que ainda carregam oxigênio, em uma molécula que também carrega CO, estão na
forma de alta afinidade: a curva se desloca para a esquerda e eles liberam menos de
seu oxigênio nos tecidos.

\textbf{22.} Pela relação de competição, a razão $[\mathrm{HbCO}]/[\mathrm{HbO_2}]$ cai em proporção a $p(\ce{O2})$:
respirar oxigênio puro, quase cinco vezes a pressão no ar, desloca o monóxido
de carbono mais depressa.

\textbf{23.} Ao dificultar a ligação linear do CO e fazer ligação de hidrogênio com o $\ce{O2}$, ela mantém $M$ na
casa das centenas, em vez do valor muito maior do \omterm{def:b3:bioinorganic:haem}{heme} livre; caso contrário, mesmo o
monóxido de carbono produzido no organismo bloquearia boa parte da hemoglobina.

\textbf{24.} Em repouso, o sangue fornece cerca de \qty{0.34}{L} de oxigênio por minuto.
\end{solution}
